% Drop-in replacement for the construction appendix.
% Requires amsmath, amssymb, amsthm, and the paper's theorem environments.
% All citation keys are already in the current bibliography or its three
% previously supplied additions. No further bibliography entries are added.
\providecommand{\R}{\mathbb R}
\providecommand{\C}{\mathbb C}
\providecommand{\Om}{\Omega}
\providecommand{\eps}{\varepsilon}
\providecommand{\dist}{\operatorname{dist}}
\providecommand{\sech}{\operatorname{sech}}
\providecommand{\norm}[1]{\ensuremath{\left\lVert #1\right\rVert}}
\providecommand{\Erep}{E_{\mathrm{rep}}}
\providecommand{\Cbd}{C_{\mathrm{bd}}}
\providecommand{\Dlam}{D_\lambda}
\providecommand{\capp}{c_{\mathrm{app}}}
\providecommand{\QUILLMainTheorem}{Theorem~\ref{thm:qi-numerics}}
\providecommand{\QUILLMainCorollary}{Corollary~\ref{cor:qi-log-width}}
% Use a dedicated reference command so a template's redefinition of \eqref
% cannot produce "equation equation ...".
\providecommand{\qeqref}[1]{(\ref{#1})}

\appendix
\numberwithin{equation}{section}
\section{Construction theory}
\label{app:construction-theory}

We prove Theorem~\ref{thm:qi-numerics-explicit}, the precise statement underlying
\QUILLMainTheorem. We first construct an accurate reference network in the
prescribed tanh dictionary, and then establish its recovery from in-domain
samples and the resulting continuous-RMS bound. Section~\ref{app:gelu-silu}
extends the representation argument to exact GELU and SiLU.

\subsection{Setup and theorem statement}
\label{sec:setup}
Let $\Om=[-1,1]$, $I_\Delta=[-1-\Delta,1+\Delta]$, and $\delta,\Delta,B>0$.
The class $\mathcal A_{\delta,B}(I_\Delta)$ consists of the functions real on
$I_\Delta$, holomorphic on
\[
 U_\delta(I_\Delta)=\{z\in\C:\dist(z,I_\Delta)<\delta\},
\]
and bounded there by $B$. We use
\[
 \norm{g}_{2,\Om}^2=\frac12\int_{-1}^1|g(x)|^2\,dx,
 \qquad \norm{g}_\infty=\sup_{x\in\Om}|g(x)|.
\]
Fix $0<\lambda\le1$ and $a_H>0$, and define
\begin{equation}
 h=\frac2N,\quad x_j=-1+jh,\quad
 R=\lceil a_H\sqrt N\rceil,\quad \gamma=\frac\lambda h,
 \quad W=N+2R+1.
 \label{eq:geometry}
\end{equation}
The tanh dictionary is
\begin{equation}
 q_c(x)=b_c+\sum_{j=-R}^{N+R}w_j\tanh(\gamma(x-x_j)),
 \qquad \norm{c}_1=|b_c|+\sum_j|w_j|.
 \label{eq:dictionary}
\end{equation}
Call this geometry admissible if
\begin{equation}
 (R+\tfrac12)h\le\Delta,\quad \gamma\delta\ge4\pi,
 \quad \lambda R\ge2\log2,\quad
 N\ge16\max\{1,a_H^2\}.
 \label{eq:admissible}
\end{equation}
At observation density $s\in\mathbb N$, use only core samples
\begin{equation}
 z_i=-1+\frac{ih}{s},\quad 0\le i\le Ns,
 \qquad M=Ns+1,\qquad
 \norm{g}_{2,X}^2=\frac1M\sum_{i=0}^{M-1}|g(z_i)|^2.
 \label{eq:samplegrid}
\end{equation}
There are $n=W+1$ readout coordinates, including the output bias. The sampled
algorithm is reference-scaled truncated SVD, as specified in
Section~\ref{sec:recovery}; it is not an unregularized least-squares minimizer.
It uses real target samples and deterministic coefficient envelopes; the
reference construction's high derivatives and complex target values are not
evaluated by the sampled algorithm. The guarantees below concern successful
runs satisfying the stated numerical conditions.

\begin{theorem}[Approximation and finite-precision recovery]
\label{thm:qi-numerics-explicit}
Let $f\in\mathcal A_{\delta,B}(I_\Delta)$ and let \qeqref{eq:geometry} be
admissible. Suppose the numerical and sampling conditions in
Section~\ref{sec:recovery} and the input-rounding conditions in
Section~\ref{sec:continuous} hold, with $M\ge n$.
Every returned network satisfies
\begin{align}
 \norm{\widehat c_W}_1
 &\le 2BK_{\rm ref}+\eps_{\rm coeff,W},
 \label{eq:reader-coeff}\\
 \norm{f-\widehat q_W}_{2,\Om}
 &\le C_f\left(e^{-\capp W}+\frac{e^{-\pi^2/\lambda}}W\right)
       +\eps_{\rm rec,W},
 \label{eq:reader-error}\\
 \capp&=\tfrac18\min\{\pi\delta,\lambda a_H^2\},
 \qquad C_f=\frac{512\pi B}{\lambda\delta}.
 \label{eq:capp}
\end{align}
Here $K_{\rm ref}$ is a width-independent bound on the sum of the reference
envelopes in \qeqref{eq:coordinate-envelope}. The nonnegative allowances
$\eps_{\rm coeff,W}$ and $\eps_{\rm rec,W}$ are defined in
\qeqref{eq:epsrec}. All input-rounding contributions
are included in these allowances; in particular, $C_f$ does not depend on $W$.
\end{theorem}

\subsection{Tanh reference network}
% Former subsection labels remain aliases of this consolidated subsection.
\label{sec:local}
\label{sec:quadrature}
\label{sec:boundary}
\label{sec:coefficients}
\label{sec:witness}

\paragraph{Local analytic representation.}
Set
\begin{equation}
 d=\frac\pi{2\gamma},\quad r=(R+\tfrac12)h,
 \quad J=[a,b]=[-1-r,1+r],\quad x_0=-1.
 \label{eq:local-geometry}
\end{equation}
The $W$ centers in \qeqref{eq:geometry} are the midpoints of the width-$h$
cells partitioning $J$, and $Wh=2+2r$. Define
\begin{equation}
 K_\gamma(t)=\frac\gamma2\sech^2(\gamma t),\qquad
 a_\gamma(z)=\frac{f(z+id)-f(z-id)}{2id}
            =\frac1{2d}\int_{-d}^d f'(z+iv)\,dv.
 \label{eq:density}
\end{equation}
Approximating a derivative by localized kernels and integrating the expansion
is a standard route to sigmoidal approximation
\cite{Costarelli2015Approximation,Buhmann_Jager_2022quasiinterpolation}.
Here the shifted-value formula gives a density adapted to a finite analytic
neighborhood.

\begin{lemma}[Local density bound]
\label{lem:density}
If $\sigma+d<\delta$, the density is holomorphic near
$J+i[-\sigma,\sigma]$, is real on $J$, and satisfies
\begin{equation}
 \sup_{t\in J,\ |y|\le\sigma}|a_\gamma(t+iy)|
 \le M_\sigma:=\frac B{\delta-\sigma-d}.
 \label{eq:density-bound}
\end{equation}
\end{lemma}
\begin{proof}
Cauchy's estimate gives $|f'(t+i(y+v))|\le B/(\delta-\sigma-d)$ for
$|v|\le d$; average in \qeqref{eq:density}. The shifted formula gives
holomorphy, and $f(\bar z)=\overline{f(z)}$ gives reality on $J$.
\end{proof}
For $x\in\Om$, put
\begin{align}
 D_x(z)&=\tfrac12[\tanh(\gamma(x-z))-\tanh(\gamma(x_0-z))],
 \qquad F_x(z)=a_\gamma(z)D_x(z), \label{eq:D-F}\\
 (\mathcal C_h f)(x)&=f(x_0)+\int_a^bF_x(t)\,dt, \label{eq:Ch}\\
 H_x(z)&=\coth(\gamma(x-z))-\coth(\gamma(x_0-z)), \label{eq:Hx}\\
 (\mathcal B_hf)(x)&=\frac1{4d}\int_{-d}^d
 [f(a+iy)H_x(a+iy)-f(b+iy)H_x(b+iy)]\,dy. \label{eq:Bh}
\end{align}

\begin{lemma}[Finite local representation]
\label{lem:finite-rep}
For $x\in\Om$,
\begin{equation}
 f(x)=\mathcal C_hf(x)+\mathcal B_hf(x),\qquad
 \norm{\mathcal B_hf}_\infty
 \le B[\coth(\gamma r)-\coth(\gamma(r+2))]
 \le\frac{2B}{e^{2\gamma r}-1}.
 \label{eq:finite-rep}
\end{equation}
\end{lemma}
\begin{proof}
The assertion at $x=x_0$ is immediate. Otherwise, $\gamma d=\pi/2$ implies
$H_x(t\pm id)=2D_x(t)$, and hence
\[
 \int_a^bF_x(t)\,dt=\frac1{4id}\left[
 \int_a^b f(t+id)H_x(t+id)\,dt-
 \int_a^b f(t-id)H_x(t-id)\,dt\right].
\]
On the positively oriented rectangle $J+i[-d,d]$, the residues of $fH_x$
at $x$ and $x_0$ are $-f(x)/\gamma$ and $f(x_0)/\gamma$.
Separating the two horizontal integrals and the vertical sides gives
\qeqref{eq:finite-rep}, since $4d\gamma=2\pi$.
For $\ell=x-x_0\in[0,2]$, the expansion
$\coth w=1+2\sum_{k\ge1}e^{-2kw}$ for $\Re w>0$ and oddness give
\begin{align*}
 |H_x(a+iy)|&\le\coth(\gamma r)-\coth(\gamma(r+\ell)),\\
 |H_x(b+iy)|&\le\coth(\gamma(r+2-\ell))-\coth(\gamma(r+2)).
\end{align*}
Each is at most $\coth(\gamma r)-\coth(\gamma(r+2))$.
Insert these bounds in \qeqref{eq:Bh} and use
$\coth t-1=2/(e^{2t}-1)$.
\end{proof}
Since $\mathcal B_h f$ depends only on the two vertical segments through
$a$ and $b$, the representation gives an explicit boundary functional.
Discretizing the integral introduces a second endpoint contribution, which
will be corrected together with this functional.

\paragraph{Midpoint discretization.}
We use the finite-interval characteristic-function argument of
\cite[Secs.~4--5]{JavedTrefethen2014}. The following formulation fixes the
orientations and retains the residues of a meromorphic integrand.
Let $b=a+m_qh$, $t_j=a+(j+\tfrac12)h$, and define
\[
 Q_hF=h\sum_{j=0}^{m_q-1}F(t_j),\qquad IF=\int_a^bF(t)\,dt,
 \qquad \zeta_0=a+h/2.
\]
For $F$ meromorphic near $[a,b]+i[-\sigma,\sigma]$, with finitely many poles
and none on the boundary or real segment, put
\begin{equation}
 \mathcal K_h(z)=\frac1{2i}\cot\frac{\pi(z-\zeta_0)}h,
 \qquad R_h^\pm=\mathcal K_h\pm\tfrac12.
 \label{eq:characteristic}
\end{equation}
Let $\mathcal P_\pm$ be its upper/lower pole sets, and set
\begin{align}
 \mathcal H_h[F]&=\int_a^b
 [R_h^-(t-i\sigma)F(t-i\sigma)-R_h^+(t+i\sigma)F(t+i\sigma)]\,dt,
 \label{eq:Hquad}\\
 \mathcal V_h[F]&=i\int_0^\sigma
 \frac{F(b+iy)-F(a+iy)-F(b-iy)+F(a-iy)}{1+e^{2\pi y/h}}\,dy,
 \label{eq:Vquad}\\
 \mathcal P_h[F]&=2\pi i\left[
 \sum_{p\in\mathcal P_+}\operatorname{Res}(R_h^+F;p)
 +\sum_{p\in\mathcal P_-}\operatorname{Res}(R_h^-F;p)\right].
 \label{eq:Pquad}
\end{align}
For a simple pole, $\operatorname{Res}(R_h^\pm F;p)
=R_h^\pm(p)\operatorname{Res}(F;p)$.

\begin{lemma}[Signed midpoint identity]
\label{lem:midpoint}
\begin{equation}
 Q_hF-IF=\mathcal H_h[F]+\mathcal V_h[F]-\mathcal P_h[F],
 \label{eq:midpoint}
\end{equation}
and, in their respective half-planes,
\begin{equation}
 |R_h^+(z)|\le\frac1{e^{2\pi\Im z/h}-1},\qquad
 |R_h^-(z)|\le\frac1{e^{2\pi|\Im z|/h}-1}.
 \label{eq:Rbounds}
\end{equation}
\end{lemma}
\begin{proof}
Apply the contour argument in \cite[Secs.~4--5]{JavedTrefethen2014} with
\qeqref{eq:characteristic}. Equivalently, the positively oriented upper
path $b\to b+i\sigma\to a+i\sigma\to a$ and lower path
$a\to a-i\sigma\to b-i\sigma\to b$ give
\[
 \int_{\Gamma_+}F=2\pi i\sum_{\mathcal P_+}\operatorname{Res}F-IF,
 \qquad
 \int_{\Gamma_-}F=2\pi i\sum_{\mathcal P_-}\operatorname{Res}F+IF.
\]
Every grid pole of $\mathcal K_h$ has residue $h/(2\pi i)$. Apply the residue
theorem to $\mathcal K_hF$ on the combined contour and substitute
$\mathcal K_h=R_h^+-1/2$ above and $R_h^-+1/2$ below. At the half-cell
endpoints, $R_h^+(a+iy)=R_h^+(b+iy)=(1+e^{2\pi y/h})^{-1}$ and the lower
values are their negatives. These identities give \qeqref{eq:midpoint}.
Writing $w=e^{2\pi i(z-\zeta_0)/h}$ gives $R_h^+=w/(w-1)$ and
$R_h^-=1/(w-1)$, proving \qeqref{eq:Rbounds}.
\end{proof}
For tanh choose
\begin{equation}
 k=\left\lceil\frac\delta{8d}\right\rceil,\qquad \sigma=2kd.
 \label{eq:sigmachoice}
\end{equation}
Admissibility implies $d\le\delta/8$ and
\begin{equation}
 \delta/4\le\sigma\le\delta/2,\quad
 \sigma+d\le5\delta/8,\quad M_\sigma\le4B/\delta.
 \label{eq:sigmaprops}
\end{equation}
Also $\gamma\sigma=k\pi$, so $D_x(t\pm i\sigma)=D_x(t)$.

\begin{lemma}[Tanh quadrature bounds]
\label{lem:tanhquad}
Uniformly for $x\in\Om$,
\begin{align}
 |\mathcal H_h[F_x]|&\le\frac{4M_\sigma}{e^{2\pi\sigma/h}-1},
 \label{eq:Hbound}\\
 |\mathcal P_h[F_x]|&\le\frac{4\pi M_\sigma}\gamma
 \sum_{\ell=0}^{k-1}\frac1{e^{(2\ell+1)\pi^2/\lambda}-1},
 \label{eq:Pbound}\\
 |\mathcal V_h[F_x]|&\le\frac{4M_\sigma h\log2}{\pi(e^{2\gamma r}-1)}.
 \label{eq:Vbound}
\end{align}
\end{lemma}
\begin{proof}
On the real axis, $D_x(t)=\int_{x_0}^xK_\gamma(v-t)\,dv\ge0$ and
$\int_\R D_x(t)\,dt=x-x_0\le2$. Together with imaginary periodicity and
\qeqref{eq:Rbounds}, this proves \qeqref{eq:Hbound}.
The pole families are
\[
 x\pm i(2\ell+1)d,\qquad x_0\pm i(2\ell+1)d,
 \qquad 0\le\ell<k.
\]
Their residues in $D_x$ have modulus $1/(2\gamma)$, and
$2\pi(2\ell+1)d/h=(2\ell+1)\pi^2/\lambda$, giving \qeqref{eq:Pbound}.
At the vertical sides, subtract the common saturated value of the two tanhs
and use $\tanh w-1=-2/(e^{2w}+1)$ for $\Re w>0$, together with oddness.
This gives $|D_x(a\pm iy)|,|D_x(b\pm iy)|\le2/(e^{2\gamma r}-1)$.
Now use \qeqref{eq:Vquad} and
$\int_0^\infty(1+e^{2\pi y/h})^{-1}\,dy=h\log2/(2\pi)$.
\end{proof}
The midpoint network is
\begin{equation}
 q_h(x)=f(x_0)+h\sum_{j=-R}^{N+R}a_\gamma(x_j)D_x(x_j).
 \label{eq:baseline-network}
\end{equation}
Its weights $w_j^{\rm base}=ha_\gamma(x_j)/2$ satisfy
\begin{equation}
 |w_j^{\rm base}|\le\frac{Bh}\delta,\qquad
 \sum_j|w_j^{\rm base}|\le\frac{2B(1+\Delta)}\delta.
 \label{eq:baseline-coeff}
\end{equation}
Subtracting \qeqref{eq:midpoint} from the finite representation gives
\begin{equation}
 f-q_h=\mathcal B_h^{\rm eff}f+\mathcal P_h[F_x]-\mathcal H_h[F_x],
 \qquad \mathcal B_h^{\rm eff}f:=\mathcal B_hf-\mathcal V_h[F_x].
 \label{eq:baseline-signed}
\end{equation}
Thus the correction must treat both the representation boundary term and the
midpoint endpoint term.

\paragraph{Boundary correction using the existing halo.}
Set
\begin{equation}
 z_-(x)=e^{-2\gamma(x-a)},\quad z_+(x)=e^{-2\gamma(b-x)},
 \quad\theta=e^{-\gamma r},\quad z_{\pm,0}=z_\pm(x_0).
 \label{eq:zpm}
\end{equation}
All these arguments lie in $(0,\theta^2]$. For $\ell\ge1$, define
\begin{align}
 \mu_\ell^-&=\frac1{2d}\int_{-d}^d f(a+iy)e^{2i\ell\gamma y}\,dy,
 &\mu_\ell^+&=\frac1{2d}\int_{-d}^d f(b+iy)e^{-2i\ell\gamma y}\,dy,
 \label{eq:mu}\\
 \nu_\ell^-&=i(-1)^\ell\int_0^\sigma
 \frac{a_\gamma(a-iy)e^{-2i\ell\gamma y}-a_\gamma(a+iy)e^{2i\ell\gamma y}}
 {1+e^{2\pi y/h}}\,dy, \label{eq:numinus}\\
 \nu_\ell^+&=i(-1)^\ell\int_0^\sigma
 \frac{a_\gamma(b-iy)e^{2i\ell\gamma y}-a_\gamma(b+iy)e^{-2i\ell\gamma y}}
 {1+e^{2\pi y/h}}\,dy. \label{eq:nuplus}
\end{align}
Write $g_\ell^\pm=\mu_\ell^\pm-\nu_\ell^\pm$,
$G^\pm(z)=\sum_{\ell\ge1}g_\ell^\pm z^\ell$, and
$\Dlam=\pi/(2\lambda)+4\log2/\pi$.

\begin{lemma}[Boundary series]
\label{lem:boundary-series}
The coefficients are real, satisfy
\begin{equation}
 |g_\ell^\pm|\le M_h:=\frac{Bh}\delta\Dlam,
 \label{eq:gh-bound}
\end{equation}
and give
\begin{equation}
 \mathcal B_h^{\rm eff}f(x)
 =G^-(z_-(x))-G^-(z_{-,0})+G^+(z_+(x))-G^+(z_{+,0}).
 \label{eq:boundary-series}
\end{equation}
\end{lemma}
\begin{proof}
Expand $\coth w=1+2\sum_{\ell\ge1}e^{-2\ell w}$ and
$\tanh w=1+2\sum_{\ell\ge1}(-1)^\ell e^{-2\ell w}$ for $\Re w>0$;
use oddness at the opposite endpoint. For example, writing
$\Delta z_\pm^\ell=z_\pm(x)^\ell-z_{\pm,0}^\ell$, the left endpoint terms
are $H_x(a+iy)=2\sum_{\ell\ge1}\Delta z_-^\ell e^{2i\ell\gamma y}$ and
$D_x(a+iy)=\sum_{\ell\ge1}(-1)^\ell\Delta z_-^\ell e^{2i\ell\gamma y}$;
the right endpoint has a minus sign and the conjugate phase. Absolute
convergence permits termwise integration, yielding \qeqref{eq:boundary-series}.
Since $\gamma d=\pi/2$,
\begin{equation}
 \int_{-d}^d e^{2i\ell\gamma y}\,dy=0\quad(\ell\ge1).
 \label{eq:moment-cancel}
\end{equation}
Subtract $f(a)$ or $f(b)$ in the corresponding $\mu$ integral. Cauchy's bound
$|f'|\le2B/\delta$ gives $|\mu_\ell^\pm|\le Bd/\delta$.
Similarly, \qeqref{eq:sigmaprops} and the integral used in the preceding
lemma give $|\nu_\ell^\pm|\le4Bh\log2/(\pi\delta)$.
These estimates imply \qeqref{eq:gh-bound}; conjugate symmetry gives reality.
\end{proof}
The cancellation makes the boundary moments $O(Bh)$ and will give an
$O(Bh)$ correction readout. The quadratic exponent in the halo size instead
comes from combining a correction degree proportional to $R$ with
$-\log\theta$ proportional to $R$.

For a series $G(z)=\sum_{\ell\ge1}g_\ell z^\ell$, $|g_\ell|\le M_G$,
use the prescribed-denominator Pad\'e-type construction
\cite{BrezinskiRedivoZaglia2015}:
\begin{equation}
 D_m(z)=\prod_{i=1}^m(1+t_i z),\quad
 P_m=[D_mG]_{\le m},\quad S_m=P_m/D_m,
 \qquad 0<t_i\le\theta^{-1}.
 \label{eq:pade}
\end{equation}
Here $[\cdot]_{\le m}$ denotes Taylor truncation at zero and the $t_i$ are
distinct. On $|z|=\theta$, $|D_mG|\le 2^mM_G/(1-\theta)$.
Cauchy's coefficient estimate, summed beyond degree $m$ on
$0\le z\le\theta^2$, and $D_m(z)\ge1$ imply
\begin{equation}
 \sup_{0\le z\le\theta^2}|G(z)-S_m(z)|
 \le \frac{M_G\theta}{(1-\theta)^2}(2\theta)^m.
 \label{eq:pade-error}
\end{equation}
Take $1\le m\le\lfloor(R+1)/2\rfloor$ and the outermost halo centers
\begin{equation}
 \xi_i^-=a+(i-\tfrac12)h,\quad
 \xi_i^+=b-(i-\tfrac12)h,\quad t_i=e^{2\lambda(i-1/2)}.
 \label{eq:outer-centers}
\end{equation}
Then $t_i\le\theta^{-1}$, and
\begin{equation}
 \frac1{1+t_i z_-(x)}=\frac{1+\tanh(\gamma(x-\xi_i^-))}{2},\qquad
 \frac1{1+t_i z_+(x)}=\frac{1-\tanh(\gamma(x-\xi_i^+))}{2}.
 \label{eq:halo-identity}
\end{equation}
Let $S_m^\pm$ be \qeqref{eq:pade} for $G^\pm$ and define
\begin{equation}
 q_{\rm bd}(x)=S_m^-(z_-(x))-S_m^-(z_{-,0})
              +S_m^+(z_+(x))-S_m^+(z_{+,0}),
 \qquad q^\star=q_h+q_{\rm bd}.
 \label{eq:qbd}
\end{equation}

\begin{proposition}[Correction in the prescribed dictionary]
\label{prop:boundary-correction}
The correction uses only the existing halo slots and satisfies
\begin{equation}
 \norm{\mathcal B_h^{\rm eff}f-q_{\rm bd}}_\infty
 \le\frac{4M_h\theta}{(1-\theta)^2}(2\theta)^m.
 \label{eq:bd-general}
\end{equation}
With $m=\lfloor(R+1)/2\rfloor$ and $\lambda R\ge2\log2$,
\begin{equation}
 \norm{\mathcal B_h^{\rm eff}f-q_{\rm bd}}_\infty
 \le16M_h e^{-\lambda R^2/4}.
 \label{eq:bd-r2}
\end{equation}
\end{proposition}
\begin{proof}
Distinct roots give $S_m=c_0+\sum_{i=1}^m c_i/(1+t_i z)$.
Use \qeqref{eq:halo-identity} in the four anchored evaluations; constants
cancel and \qeqref{eq:pade-error} proves \qeqref{eq:bd-general}.
Now $m\ge R/2$, $\theta\le e^{-\lambda R}\le1/4$, and
$(2\theta)^m\le e^{-\lambda R^2/4}$, while
$4\theta/(1-\theta)^2\le16$.
\end{proof}

\paragraph{Width-independent readout control.}
The partial-fraction coefficients form an inverse-Vandermonde map
\cite{Gautschi1962}. For the nodes in \qeqref{eq:outer-centers}, put
\begin{equation}
 \zeta=e^{-2\lambda},\quad r_i=t_i^{-1}=\zeta^{i-1/2},\quad
 p_k=\prod_{j=1}^k(1-\zeta^j),\quad p_0=1.
 \label{eq:zeta-p}
\end{equation}
The coefficients in $S_m=c_0+\sum_i c_i/(1+t_i z)$ satisfy
\begin{equation}
 c_i=(-1)^i\frac{\zeta^{[i(i+1)-1]/2}}{p_{i-1}p_{m-i}}
 \sum_{\ell=1}^m e_{\ell-1}(r_1,\ldots,\widehat r_i,\ldots,r_m)g_\ell.
 \label{eq:scaled-pf}
\end{equation}
Here $e_k$ is the elementary symmetric polynomial and a hat omits an entry.
For clarity, the specialization of the classical coefficient formula is
\[
 c_i=\frac{P_m(-r_i)}{\prod_{j\ne i}(1-t_j/t_i)},\qquad
 P_m(-r_i)=(-r_i)^m\sum_{\ell=1}^m g_\ell
 [z^{m-\ell}]\prod_{j\ne i}(1+t_jz).
\]
The second equality follows by telescoping the factor $1+t_i z$ in the
truncated numerator. Factoring the geometric node differences in the
first equality gives \qeqref{eq:scaled-pf}.
Define
\begin{align}
 p_\infty&=\prod_{j\ge1}(1-\zeta^j),\qquad
 L_+=\prod_{j\ge1}(1+\zeta^{j-1/2}), \label{eq:pinfty}\\
 \Cbd(\lambda)&=\sum_{i\ge1}
 \frac{\zeta^{[i(i+1)-1]/2}}{p_{i-1}p_\infty}
 \prod_{\substack{j\ge1\\j\ne i}}(1+\zeta^{j-1/2}).
 \label{eq:Cbd}
\end{align}

\begin{proposition}[Boundary coefficient bound]
\label{prop:Cbd}
For fixed $\lambda>0$, $\Cbd(\lambda)<\infty$ and
\begin{equation}
 \sum_{i=1}^m|c_i|\le\Cbd(\lambda)\sup_{\ell\ge1}|g_\ell|,
 \qquad
 \Cbd(\lambda)\le\frac{\sqrt\pi e^\lambda}{2\sqrt\lambda}
                  \exp\frac{5\pi^2}{24\lambda}.
 \label{eq:Cbd-bound}
\end{equation}
\end{proposition}
\begin{proof}
In \qeqref{eq:scaled-pf}, the symmetric coefficients are nonnegative and
sum to $\prod_{j\ne i}(1+r_j)$; also $p_{m-i}\ge p_\infty$.
This proves the first bound. Since $\sum_j\zeta^j<\infty$, $p_\infty>0$
and $L_+<\infty$, and the remaining $i$-factor decays as
$e^{-\lambda i(i+1)}$. For the explicit bound, use $p_{i-1}\ge p_\infty$ and
\begin{align*}
 \log p_\infty^{-1}
 &=\sum_{k\ge1}\frac1{k(e^{2\lambda k}-1)}\le\frac{\pi^2}{12\lambda},\\
 \log L_+&\le\frac1{2\lambda}\int_0^\infty\log(1+e^{-t})\,dt
          =\frac{\pi^2}{24\lambda},\\
 \sum_{i\ge1}e^{-\lambda i(i+1)}
 &\le\int_0^\infty e^{-\lambda t^2}\,dt=\frac{\sqrt\pi}{2\sqrt\lambda}.
\end{align*}
The middle inequality is the midpoint bound for the convex integrand.
\end{proof}
Each boundary contributes half of its partial-fraction absolute sum to the
tanh readout. Thus
\begin{equation}
 \sum_j|w_j^\star|\le\frac{2B(1+\Delta)}\delta+
                 \Cbd(\lambda)\frac{Bh}\delta\Dlam.
 \label{eq:weight-sum}
\end{equation}
The anchored bias is at most $B+\sum_j|w_j^\star|$, so $h\le1/8$ gives
\begin{equation}
 \norm{c^\star}_1\le K B,\qquad
 K=1+\frac{4(1+\Delta)}\delta+\frac{\Cbd(\lambda)\Dlam}{4\delta}.
 \label{eq:K}
\end{equation}

\paragraph{The analytic tanh witness.}
\begin{theorem}[Boundary-corrected analytic witness]
\label{thm:witness}
For every admissible geometry,
\begin{equation}
 \norm{f-q^\star}_\infty\le
 \frac{128\pi B}{\lambda\delta}
 \left[e^{-\capp W}+\frac{e^{-\pi^2/\lambda}}W\right]
 =:\Erep(W,\lambda),\qquad \norm{c^\star}_1\le KB.
 \label{eq:Erep}
\end{equation}
In particular, $W=N+O(\sqrt N)$ and $\gamma=\Theta(W)$ at fixed $\lambda$.
\end{theorem}
\begin{proof}
The signed decomposition is
\begin{equation}
 f-q^\star=(\mathcal B_h^{\rm eff}f-q_{\rm bd})
            +\mathcal P_h[F_x]-\mathcal H_h[F_x].
 \label{eq:corrected-decomp}
\end{equation}
For $A\ge\pi^2$,
\[
 \sum_{\ell\ge0}\frac1{e^{(2\ell+1)A}-1}
 \le\frac{e^{-A}}{(1-e^{-A})(1-e^{-2A})}\le2e^{-A}.
\]
Consequently, Lemma~\ref{lem:tanhquad} and
Proposition~\ref{prop:boundary-correction} give the explicit three bounds
\begin{align*}
 E_{\rm hor}&\le\frac{32B}\delta e^{-\pi\delta N/4},\\
 E_{\rm bd}&\le\frac{32B\Dlam}{\delta N}e^{-\lambda a_H^2N/4},\\
 E_{\rm rep\!l}&\le\frac{64\pi B}{\lambda\delta N}e^{-\pi^2/\lambda}.
\end{align*}
Admissibility implies $N\le W\le2N$ and
$\lambda(1+\Dlam)<4\pi$. With
$\kappa=\tfrac14\min\{\pi\delta,\lambda a_H^2\}$, the first two terms
sum to at most $(128\pi B/(\lambda\delta))e^{-\kappa W/2}$;
the third is at most $(128\pi B/(\lambda\delta W))e^{-\pi^2/\lambda}$.
Equation \qeqref{eq:K} gives the readout bound.
\end{proof}

\subsection{Recovery and continuous error}
\label{sec:recovery}

\subsubsection{Common dictionary assumptions}
The recovery argument below is conditional on the following properties of a
prescribed dictionary. They hold for the tanh construction just proved; the
GELU/SiLU construction in Section~\ref{app:gelu-silu} will verify them before
applying the same bounds.
Assume $M\ge n$, where $n$ is the total readout dimension, including the bias
and any affine-correction units. Write the native features as $\phi_j$, with
$\phi_b=1$, and set
\begin{equation}
 \varphi_j=\phi_j/m_j,\quad
 \mathsf M_\phi=\operatorname{diag}(m_j),\quad v=\mathsf M_\phi c,
 \quad q_c=\sum_jv_j\varphi_j,
 \qquad m_b=1.
 \label{eq:normalized-dictionary}
\end{equation}
Assume $m_j\ge1$ and $|\varphi_j(x)|\le1$ on $\Om$. Assume also that a
reference network satisfies $\norm{f-q^\star}_\infty\le E_W^{\rm an}$ and
has positive deterministic coordinate envelopes
\begin{equation}
 |v_j^\star|\le B\alpha_j,\qquad
 K_W^{\rm ref}:=\sum_j\alpha_j\le K_{\rm ref}.
 \label{eq:coordinate-envelope}
\end{equation}
Here $K_{\rm ref}$ is independent of width for the fixed design.

For the continuous-error conclusions, use the core-cell Bernstein ellipses
and impose the following sampling conditions.
Fix $\varrho>1$ and $p\in\mathbb N$, $p\ge1$, and assume
\begin{equation}
 \lambda(\varrho-\varrho^{-1})\le\pi,\qquad
 \frac h4(\varrho+\varrho^{-1})<\delta,\qquad s\ge2(p+1)^2.
 \label{eq:sampling-conditions}
\end{equation}
Assume that the target and normalized features are holomorphic on a
neighborhood of every closed cell ellipse, with $|f(z)|\le B$ and
$|\varphi_j(z)|\le C_{\rm ell}$ there for a width-independent constant
$C_{\rm ell}$. The numerical comparison and export conditions
\qeqref{eq:comparison-objects}--\qeqref{eq:nearby-problem}, cutoff condition
\qeqref{eq:eta}, and input-rounding and activation-evaluation assumptions
below are additional conditions on a returned run; they do not follow from
these dictionary properties.

\paragraph{Verification for tanh.}
Take $m_j=1$ and $E_W^{\rm an}=\Erep(W,\lambda)$ from
Theorem~\ref{thm:witness}. The real-axis feature bound is immediate.
Put $m=\lfloor(R+1)/2\rfloor$ and
\begin{equation}
 L_{i,m}=\frac{\zeta^{[i(i+1)-1]/2}}{p_{i-1}p_{m-i}}
             \prod_{\substack{1\le j\le m\\j\ne i}}(1+\zeta^{j-1/2}).
 \label{eq:Lim}
\end{equation}
Start each grid coordinate at $\alpha_j=h/[2(\delta-d)]$, and add
$h\Dlam L_{i,m}/(2\delta)$ at the $i$th corrected slot on either side.
Take $\alpha_b=1+\sum_{j\ne b}\alpha_j$. The coefficient formulas above give
\qeqref{eq:coordinate-envelope}; one may use $K_{\rm ref}=K$ from
\qeqref{eq:K}, because $\sum_iL_{i,m}\le\Cbd(\lambda)$.
The closed cell ellipses lie in the target tube. In preactivation coordinates
their vertical semiaxis is at most $\pi/4$. For tanh,
\[
 |\tanh(u+iv)|^2=\frac{\sinh^2u+\sin^2v}{\sinh^2u+\cos^2v}\le1
 \qquad(|v|\le\pi/4).
\]
Thus the complex feature hypothesis holds with $C_{\rm ell}=1$.

\subsubsection{Reference scaling}
Define
\begin{equation}
 D=\operatorname{diag}(\sqrt{\alpha_j}),\quad
 L=\sqrt{\sum_j\alpha_j},\quad S=BL,\quad
 A_{ij}=\varphi_j(z_i),\quad
 F=\frac{AD}{\sqrt M},\quad y_X=\frac{(f(z_i))_{i=0}^{M-1}}{\sqrt M}.
 \label{eq:scaling}
\end{equation}
These are exact reference quantities; rounding of their stored versions is
included in the formation bounds below. Cauchy--Schwarz gives
\begin{equation}
 a^\star=D^{-1}v^\star,\qquad
 \norm{a^\star}_2\le S,\quad
 \norm{Fa^\star-y_X}_2\le E_W^{\rm an},\quad
 \norm{Da}_1\le L\norm{a}_2.
 \label{eq:scaled-witness}
\end{equation}
Indeed, $\sum_j|v_j^\star|^2/\alpha_j\le B^2\sum_j\alpha_j$;
the residual bound follows from uniform approximation.

\subsubsection{Numerical conditions and comparison problem}
Use round-to-nearest arithmetic with unit roundoff $u$, gradual underflow,
and no overflow or invalid operations in a returned run. Let $\mu$ be a
valid absolute rounding allowance (the smallest positive subnormal is
sufficient), so $|\operatorname{RN}(t)-t|\le u|t|+\mu$.
All error bounds are either standard a priori bounds or validated upper
bounds with outward rounding. Failed range or denominator checks do not
produce a certified return. Standard product, summation, and matrix-operation
estimates are taken from \cite[Chs.~2--4]{Higham2002}; SVD factorization
and orthogonality bounds are supplied by \cite[Ch.~4]{LAPACK1999}.

For later use, a length-$q$ pairwise dot product with rounded products obeys
\begin{equation}
 |\operatorname{pdot}(v,w)-v^Tw|
 \le\Gamma_q^{\rm pw}\sum_i|v_iw_i|+\beta_q^{\rm pw},\quad
 \Gamma_q^{\rm pw}=\frac{k_qu}{1-k_qu},\quad
 \beta_q^{\rm pw}=\frac{(2q-1)\mu}{1-k_qu},\quad
 k_q=1+\lceil\log_2q\rceil.
 \label{eq:pairwise-bound}
\end{equation}
Assume $k_qu<1$. The relative bound is the usual pairwise bound; the additive
term counts at most $2q-1$ local absolute errors and their amplification
along the tree. In particular $\Gamma_q^{\rm pw}=O(u\log q)$.

The algorithm assembles $\widetilde F,\widetilde y$, computes a thin SVD
$(\widehat U,\widehat\Sigma,\widehat V)$ with nonnegative diagonal
$\widehat\Sigma$, forms $p_i=\operatorname{pdot}(\widehat U_{:i},\widetilde y)$,
and retains exactly the singular values $\widehat\sigma_i>\tau$. It sets
\begin{equation}
 \widehat z_i=\begin{cases}
 \operatorname{RN}(p_i/\widehat\sigma_i),&\widehat\sigma_i>\tau,\\
 0,&\widehat\sigma_i\le\tau,
 \end{cases}
 \label{eq:computed-svd-coordinates}
\end{equation}
and exports the native readout by forming
$D\widehat V\widehat z$ and dividing coordinate $j$ by $m_j$, all in the
working arithmetic. The bias is included in this same solve.

We require bounds for the exact comparison quantities
\begin{equation}
 F_0=U_0\widehat\Sigma V_0^T,\quad U_0^TU_0=V_0^TV_0=I,
 \qquad
 a_\tau=F_{0,\tau}^{\dagger}y_0=V_0\widehat z,
 \label{eq:comparison-objects}
\end{equation}
with
\begin{equation}
 \norm{F-F_0}_2\le a_W,\quad \norm{y_X-y_0}_2\le b_W,\quad
 \norm{\mathsf M_\phi\widehat c-Da_\tau}_1\le e_{c,W}.
 \label{eq:nearby-problem}
\end{equation}
The matrix and label formation bounds include normalization, evaluation at
stored sample points, and any hidden-geometry storage error. Thus the analytic
basis remains the prescribed basis, rather than an implicitly perturbed one.

Here is a sufficient realization of \qeqref{eq:comparison-objects}--\qeqref{eq:nearby-problem}
that also makes the dependence on the cutoff explicit. Suppose
\begin{align*}
 \norm{F-\widetilde F}_2&\le e_F,&
 \norm{y_X-\widetilde y}_2&\le e_y,&
 \norm{\widetilde F-\widehat U\widehat\Sigma\widehat V^T}_2&\le r_F,\\
 \norm{\widehat U^T\widehat U-I}_2&\le o_U<1,&
 \norm{\widehat V^T\widehat V-I}_2&\le o_V<1.
\end{align*}
Put $\delta_U=1-\sqrt{1-o_U}$, $\delta_V=1-\sqrt{1-o_V}$,
$U_+=\sqrt{1+o_U}$, $V_+=\sqrt{1+o_V}$, and $Y=\norm{\widetilde y}_2$.
Standard bounds or enclosed residual products supply these numbers; spectral
norms may be bounded by Frobenius norms. Then one may take
\begin{align}
 a_W&=e_F+r_F+\norm{\widehat\Sigma}_2(\delta_U V_++\delta_V),\nonumber\\
 e_p&=\sqrt n\,[\Gamma_M^{\rm pw}U_+Y+\beta_M^{\rm pw}],\nonumber\\
 b_W&=e_y+\delta_UY+e_p+u\norm p_2+\mu\norm{\widehat\sigma}_2.
 \label{eq:comparison-bounds}
\end{align}
To verify the interface, take the polar factors
$U_0=\widehat U(\widehat U^T\widehat U)^{-1/2}$ and similarly $V_0$.
Their singular values give
$\norm{\widehat U-U_0}_2\le\delta_U$ and
$\norm{\widehat V-V_0}_2\le\delta_V$, proving the bound for $a_W$.
If $J_\tau=\{i:\widehat\sigma_i>\tau\}$, set
\[
 y_0=\widetilde y+U_{0,J_\tau}
 (\widehat\Sigma_{J_\tau}\widehat z_{J_\tau}
       -U_{0,J_\tau}^T\widetilde y).
\]
This makes $F_{0,\tau}^{\dagger}y_0=V_0\widehat z$ exactly. Projection
rounding contributes $e_p$, and division gives
$|\widehat\sigma_i\widehat z_i-p_i|\le u|p_i|+\mu\widehat\sigma_i$ on
retained coordinates. These bounds prove the formula for $b_W$.
In particular, $a_W$ and $b_W$ are available before choosing $\tau$; the
bound for division does not contain $1/\tau$.
The polar factors and $y_0$ are analysis objects, not computed intermediates.

Finally let $e_{\rm exp}$ be the standard forward-error bound for the
specified native export, in the form
\[
 \norm{\mathsf M_\phi\widehat c-D\widehat V\widehat z}_1\le e_{\rm exp}.
\]
It includes storage of $D$, row dot products, products, and division by $m_j$.
It can also be obtained by enclosing each displayed coordinate difference
using \qeqref{eq:pairwise-bound} and the scalar rounding model. Once the
bound $\norm{\widehat z}_2\le2S$ below is established, the choice
\begin{equation}
 e_{c,W}=e_{\rm exp}+2LS\delta_V
 \label{eq:coefficient-export-bound}
\end{equation}
proves the last inequality of \qeqref{eq:nearby-problem}. This construction
of the numerical interface uses only standard error bounds and returned
finite arrays; it does not assume recovery of the analytic coefficients.

\subsubsection{Cutoff and recovery bound}
Let $\overline E_W\ge E_W^{\rm an}$ be the approximation bound available to
the algorithm, and retain its formation allowance
$\eps_{E,W}:=\overline E_W-E_W^{\rm an}\ge0$.
Choose a positive threshold with
\begin{equation}
 \eta_W=\overline E_W+a_WS+b_W,\qquad
 \tau S\ge\eta_W,\qquad 0\le\tau S-\eta_W\le\rho_W.
 \label{eq:eta}
\end{equation}
An outward upper bound for $\eta_W/S$ provides such a threshold; $\rho_W$
includes this excess and its rounding. No condition on singular-value gaps
is needed. The truncated-SVD coefficient comparison in
\cite[Theorem~3.8, first inequality]{AdcockHuybrechs2020} gives
\begin{equation}
 \norm{a_\tau}_2\le S+\eta_W/\tau\le2S,
 \qquad \norm{F_0a_\tau-y_0}_2\le\eta_W+\tau S.
 \label{eq:tsvd-standard}
\end{equation}
For the second inequality, the discarded singular values are at most $\tau$;
apply this to the comparison vector $a^\star$. This is the usual
thresholded least-squares estimate \cite{Hansen1987TSVD}.

\begin{proposition}[Recovered coefficient and sample bounds]
\label{prop:svd-recovery}
Define
\begin{equation}
 T=2LS+e_{c,W},\qquad
 R_X=2\overline E_W+4a_WS+3b_W+\rho_W+e_{c,W}.
 \label{eq:T-RX}
\end{equation}
The returned native readout satisfies
\begin{equation}
 \norm{\widehat c}_1\le\norm{\mathsf M_\phi\widehat c}_1\le T,
 \qquad \norm{f-q_{\widehat c}}_{2,X}\le R_X.
 \label{eq:sampled-recovery}
\end{equation}
\end{proposition}
\begin{proof}
The reference residual on $(F_0,y_0)$ is at most $\eta_W$, so
\qeqref{eq:tsvd-standard} applies. Since $a_\tau=V_0\widehat z$,
$\norm{\widehat z}_2=\norm{a_\tau}_2\le2S$, which justifies
\qeqref{eq:coefficient-export-bound}. Returning to $F,y_X$ gives
\[
 \norm{Fa_\tau-y_X}_2
 \le\eta_W+\tau S+2a_WS+b_W
 \le2\overline E_W+4a_WS+3b_W+\rho_W.
\]
Now use \qeqref{eq:scaled-witness}, coefficient export, and
$|\varphi_j|\le1$. Finally, $m_j\ge1$ bounds the native norm by the normalized
norm. The reference part of $T$ is $2BL^2\le2BK_{\rm ref}$.
\end{proof}
At fixed observation density, the pairwise projection contribution to
$e_p$ alone is $O(u\sqrt W\log W)$ times $U_+Y$.

\subsubsection{From sampled error to continuous error}
\label{sec:continuous}
Let $q=q_{\widehat c}$ be the exact function of the returned readout in the
prescribed basis, and put $g=f-q$. For each core cell $I=[x_k,x_{k+1}]$, define
\[
 \norm v_{2,I}^2=\frac1h\int_I|v|^2,\qquad
 \norm v_{\mathrm{trap},I}^2=
 \frac1s\left[\frac{|v(x_k)|^2+|v(x_{k+1})|^2}2
          +\sum_{j=1}^{s-1}|v(x_k+jh/s)|^2\right].
\]
Under the common dictionary assumptions, $g$ is holomorphic on the cell
ellipses and the recovered normalized coefficient bound gives
\begin{equation}
 |g(z)|\le\mathcal M:=B+C_{\rm ell}T
 \quad\text{on every cell ellipse}.
 \label{eq:gellipse}
\end{equation}
The analytic Chebyshev estimate \cite[Ch.~8]{trefethen2019atap} gives a
degree-$p$ polynomial $P_I$ such that
\begin{equation}
 \norm{g-P_I}_{\infty,I}\le e_p^{\rm an}
 :=\frac{2\mathcal M}{\varrho-1}\varrho^{-p}.
 \label{eq:ep}
\end{equation}
The local polynomial inverse estimates used below are standard; see
\cite[Secs.~5.4.1 and 5.4.4]{Canuto2006Spectral}. We record their constants
and the specialization to the present equispaced grid.

\begin{lemma}[Local polynomial sampling]
\label{lem:poly-sampling}
For every polynomial $P$ of degree at most $p$ on a cell of length $h$,
\begin{align}
 \norm{P'}_{2,I}&\le\frac{(p+1)\sqrt{p(p+2)}}h\norm P_{2,I},
 \label{eq:legendre-inverse}\\
 \norm P_{\infty,I}&\le(p+1)\norm P_{2,I},\qquad
 \norm{P'}_{\infty,I}\le\frac{2p^2(p+1)}h\norm P_{2,I}.
 \label{eq:linfty-inverse}
\end{align}
Under \qeqref{eq:sampling-conditions},
\begin{equation}
 \norm P_{2,I}\le\sqrt2\norm P_{\mathrm{trap},I},\qquad
 \norm P_{\infty,I}\le2\max_{x\in X\cap I}|P(x)|.
 \label{eq:P-trap}
\end{equation}
\end{lemma}
\begin{proof}
For completeness, the constants in the cited inverse estimates can be
obtained from the normalized Legendre basis
$\sqrt{2j+1}\,P_j$ on $[-1,1]$:
$|P_j|\le1$, $\int_{-1}^1(P_j')^2=j(j+1)$, and
$\sum_{j=1}^p(2j+1)j(j+1)/2=p(p+1)^2(p+2)/4$.
Cauchy--Schwarz and rescaling give \qeqref{eq:legendre-inverse} and the first
inequality in \qeqref{eq:linfty-inverse}; Markov's inequality gives the second.
For the equispaced grid, the elementary trapezoidal error for $|P|^2$ gives
\[
 \left|\norm P_{\mathrm{trap},I}^2-\norm P_{2,I}^2\right|
 \le\frac hs\norm P_{2,I}\norm{P'}_{2,I}
 \le\frac{(p+1)\sqrt{p(p+2)}}s\norm P_{2,I}^2
 \le\tfrac12\norm P_{2,I}^2.
\]
This proves the first part of \qeqref{eq:P-trap}. For the uniform norm,
choose the nearest grid point to a maximizer, at distance at most $h/(2s)$,
and use Markov's inequality to get
$\norm P_\infty\le\max_X|P|+(p^2/s)\norm P_\infty$.
\end{proof}

\begin{lemma}[Continuous error transfer]
\label{lem:rms-transfer}
With $c_M=\sqrt{2M/(M-1)}\le2$,
\begin{align}
 \norm g_{2,\Om}&\le c_M\norm g_{2,X}+(1+\sqrt2)e_p^{\rm an},
 \label{eq:rms-transfer}\\
 \norm g_\infty&\le2\max_{x\in X}|g(x)|+3e_p^{\rm an}.
 \label{eq:uniform-transfer}
\end{align}
\end{lemma}
\begin{proof}
On each cell, \qeqref{eq:ep} and \qeqref{eq:P-trap} give
$\norm g_{2,I}\le\sqrt2\norm g_{\mathrm{trap},I}+(1+\sqrt2)e_p^{\rm an}$.
Combine cells by Minkowski. The global trapezoidal norm squared is
\[
 \frac1{M-1}\left[\frac{|g(z_0)|^2+|g(z_{M-1})|^2}2
                       +\sum_{i=1}^{M-2}|g(z_i)|^2\right]
 \le\frac M{M-1}\norm g_{2,X}^2.
\]
For the uniform inequality, apply the second part of \qeqref{eq:P-trap} to
$P_I$, and add the three approximation-error terms.
\end{proof}

\paragraph{Input rounding and evaluation.}
Let $Q(x)$ be round-to-nearest in the working format, clamped to $\Om$, and
let $\xi\ge\sup_{x\in\Om}|Q(x)-x|$. For the RMS rounding estimate assume
that the core-cell endpoints are exactly representable. Monotonicity of
round-to-nearest then gives $Q(I)\subseteq I$ for each core cell. For example, $N=2^k$ satisfies this condition when the required grid
significands, exponents, and integer indices are exactly representable.
Set
\begin{equation}
 \chi_W=\frac{2p^2(p+1)\xi}h.
 \label{eq:chi}
\end{equation}
The bound on $P_I'$ in \qeqref{eq:linfty-inverse} implies
$\norm{P_I\circ Q}_{2,I}\le(1+\chi_W)\norm{P_I}_{2,I}$.
Repeating the proof of Lemma~\ref{lem:rms-transfer} therefore gives
\begin{equation}
 \norm{g\circ Q}_{2,\Om}
 \le(1+\chi_W)c_M R_X+
       [1+\sqrt2(1+\chi_W)]e_p^{\rm an}.
 \label{eq:gQ}
\end{equation}
Also $\norm{f-f\circ Q}_{2,\Om}\le B\xi/\delta$ by Cauchy's estimate.
We collect the rounding terms separately:
\begin{equation}
 E_{\rm in,W}=\frac{B\xi}\delta+
        \chi_W(c_M R_X+\sqrt2 e_p^{\rm an}).
 \label{eq:input-round-target}
\end{equation}

Assume the activation evaluator, including stored affine arguments, satisfies
\[
 |\widehat\phi_j(t)-\phi_j(t)|\le m_j\eps_{\phi,W}
 \quad(t\in\Om\text{ representable}).
\]
This is a uniform error contract for the specified tanh, exact-GELU, or SiLU
routine, including any geometry-storage perturbation; it is not an assumption
of an exact transcendental evaluation. Use rounded products and pairwise
summation for the native readout. Since
$\sum_j|\widehat c_j\phi_j(t)|\le T$ and
$\sum_j|\widehat c_j\widehat\phi_j(t)|\le(1+\eps_{\phi,W})T$,
\qeqref{eq:pairwise-bound} gives
\begin{equation}
 E_{\rm ev,W}:=
 [\eps_{\phi,W}+\Gamma_n^{\rm pw}(1+\eps_{\phi,W})]T+\beta_n^{\rm pw},
 \qquad
 |\widehat q_W(x)-q(Q(x))|\le E_{\rm ev,W}.
 \label{eq:evaluation}
\end{equation}
The mixed term $\Gamma_n^{\rm pw}\eps_{\phi,W}$ is retained.

\begin{theorem}[Computed continuous-error bounds]
\label{thm:computed}
Under the preceding analytic and numerical conditions,
\begin{equation}
 \norm{\widehat c_W}_1\le T,\qquad
 \norm{f-\widehat q_W}_{2,\Om}
 \le c_M R_X+(1+\sqrt2)e_p^{\rm an}+E_{\rm in,W}+E_{\rm ev,W}.
 \label{eq:computed-rms}
\end{equation}
If $r_{\infty,X}$ is a validated upper bound on
$\max_{x\in X}|f(x)-q(x)|$, then also
\begin{equation}
 \norm{f-\widehat q_W}_\infty
 \le2r_{\infty,X}+3e_p^{\rm an}+B\xi/\delta+E_{\rm ev,W}.
 \label{eq:computed-uniform}
\end{equation}
The uniform statement requires only $Q(\Om)\subseteq\Om$, not cellwise
preservation. It is a maximum-residual certificate, not a replacement of the
sampled RMS by the same numerical value.
\end{theorem}
\begin{proof}
The coefficient bound is Proposition~\ref{prop:svd-recovery}.
Decompose $f-\widehat q=(f-f\circ Q)+(g\circ Q)+(q\circ Q-\widehat q)$,
and apply \qeqref{eq:gQ} and \qeqref{eq:evaluation} for RMS.
For the uniform norm, use \qeqref{eq:uniform-transfer} at $Q(x)$.
\end{proof}

\subsubsection{Tanh theorem and width dependence}
\label{sec:assemble}
Put
\begin{equation}
 A_{\rm rep}=\frac{128\pi}{\lambda\delta},\qquad
 \Delta_{\rm rec,W}=2\eps_{E,W}+4a_WS+3b_W+\rho_W+e_{c,W}.
 \label{eq:Arep}
\end{equation}
For tanh, $R_X=2\Erep(W,\lambda)+\Delta_{\rm rec,W}$.
Using $c_M\le2$ in the analytic term of \qeqref{eq:computed-rms} gives
\begin{equation}
 \norm{f-\widehat q_W}_{2,\Om}
 \le4A_{\rm rep}B\left(e^{-\capp W}+\frac{e^{-\pi^2/\lambda}}W\right)
       +\eps_{\rm rec,W},
 \label{eq:final-decomp}
\end{equation}
where
\begin{equation}
 \eps_{\rm rec,W}=c_M\Delta_{\rm rec,W}
 +(1+\sqrt2)e_p^{\rm an}+E_{\rm in,W}+E_{\rm ev,W},
 \qquad \eps_{\rm coeff,W}=e_{c,W}.
 \label{eq:epsrec}
\end{equation}
These quantities are nonnegative explicit bounds; $E_{\rm in,W}$ contains
all factors involving $\xi/h$. The reference coefficient bound is
$2LS=2B K_W^{\rm ref}\le2BK_{\rm ref}$, which completes the proof of
Theorem~\ref{thm:qi-numerics-explicit}.

\begin{corollary}[Logarithmic sufficient width above a fixed allowance]
\label{cor:logwidth}
Fix the target class, bandwidth, activation, and design. Suppose the preceding
bounds have the form
\[
 \norm{f-\widehat q_W}_{2,\Om}\le A e^{-cW}+\tau_W
\]
with $A,c>0$ independent of width, and $\tau_W\le\tau$ throughout a specified
certified range $\mathcal W$ of admissible widths. For any $\eps>\tau$, every
$W\in\mathcal W$ satisfying
\begin{equation}
 W\ge\frac1c\log\frac A{\eps-\tau}
 \label{eq:log-width}
\end{equation}
has error at most $\eps$. For tanh one may take $A=4A_{\rm rep}B$,
$c=\capp$, and
$\tau_W=4A_{\rm rep}B W^{-1}e^{-\pi^2/\lambda}+\eps_{\rm rec,W}$.
\end{corollary}
\begin{proof}
The error is bounded by $Ae^{-cW}+\tau$; solve
$Ae^{-cW}\le\eps-\tau$. Existence of an admissible width meeting this
inequality inside the certified range remains part of the premise.
\end{proof}
For $\eps\ge2\tau$ the sufficient width is
$c^{-1}\log(2A/\eps)$, subject to the same range restriction. This is the
fixed-design statement underlying \QUILLMainCorollary. If $\lambda$ is
instead varied with $\eps$, the dependence of $A,c$, admissibility, and the
allowances on $\lambda$ must be retained; this corollary alone makes no joint
$\eps\to0$ assertion with varying design.

\subsection{Extensions to exact GELU and SiLU}
\label{app:gelu-silu}
For these activations we approximate $f''$ by a polynomial, apply a finite
differential inverse of the kernel convolution, and discretize by
Lemma~\ref{lem:midpoint}. We then verify the dictionary assumptions of
Section~\ref{sec:recovery} and apply its recovery and continuous-error bounds.

\subsubsection{Domains and polynomial approximation}
Choose $1<L_*<1+\Delta$ and let $J_*=[-L_*,L_*]$. Denote by
$\mathcal E_\rho(J_*)$ the closed filled Bernstein ellipse with semiaxes
$L_*(\rho+\rho^{-1})/2$ and $L_*(\rho-\rho^{-1})/2$.
Fix $1<\rho_1<\rho_2$, $\delta_*>0$, and $s_*>0$ such that
\begin{equation}
 \mathcal E_{\rho_2}(J_*)+\{z:|z|\le\delta_*\}
 \subset U_\delta(I_\Delta),\qquad
 \mathcal D:=J_*+i[-s_*,s_*]\Subset\operatorname{int}\mathcal E_{\rho_1}(J_*).
 \label{eq:psi-domains}
\end{equation}
Such choices exist by taking the ellipses and $s_*,\delta_*$ sufficiently
small. Set
\begin{equation}
 d_*:=\dist(\mathcal D,\partial\mathcal E_{\rho_1}(J_*))>0,
 \quad d_{\rm ext}=L_*-1,\quad C_*:=2/d_{\rm ext}.
 \label{eq:psi-margins}
\end{equation}
Cauchy's estimate bounds $f''$ on $\mathcal E_{\rho_2}(J_*)$ by
$2B/\delta_*^2$. The analytic Chebyshev estimates of
\cite[Ch.~8]{trefethen2019atap}, applied to its degree-$\nu$ truncation,
give a real polynomial $p_\nu$ with
\begin{equation}
 \norm{f''-p_\nu}_{L^\infty(J_*)}\le A_{\rm pol}B\rho_2^{-\nu},\qquad
 \sup_{\mathcal E_{\rho_1}(J_*)}|p_\nu|\le A_{\rm pol}B,
 \label{eq:psi-poly}
\end{equation}
where, for example, one may take
\begin{equation}
 A_{\rm pol}=\frac2{\delta_*^2}
 \max\left\{\frac2{\rho_2-1},\ 1+\frac{2\rho_1}{\rho_2-\rho_1}\right\}.
 \label{eq:Apol}
\end{equation}
The second inequality follows by summing the Chebyshev coefficient bound on
the smaller ellipse. We reserve $n$ for the number of readout coordinates
and use $\nu=\lfloor\eta N\rfloor$ for the polynomial degree.

\subsubsection{Finite polynomial deconvolution}
For a real even kernel $K$ with an exponential absolute moment and integral one,
write
\begin{equation}
 M_K(z)=\int_\R K(t)e^{zt}\,dt,\qquad
 M_K(z)^{-1}=\sum_{\ell\ge0}b_{2\ell}z^{2\ell}.
 \label{eq:psi-mgf}
\end{equation}
Suppose the inverse is holomorphic on $|z|\le r_*$ and bounded there by
$C_{\rm inv}$. Define
\begin{equation}
 a_{\gamma,\nu}
 =\sum_{\ell=0}^{\lfloor\nu/2\rfloor}
 b_{2\ell}\gamma^{-2\ell}p_\nu^{(2\ell)},\qquad
 A_{\rm den}:=\tfrac43 C_{\rm inv}A_{\rm pol}.
 \label{eq:psi-deconv}
\end{equation}

\begin{lemma}[Polynomial deconvolution]
\label{lem:psi-deconv}\label{lem:psi-density}
The density is a real polynomial. With $K_\gamma(t)=\gamma K(\gamma t)$,
\begin{equation}
 \int_\R a_{\gamma,\nu}(t)K_\gamma(x-t)\,dt=p_\nu(x).
 \label{eq:psi-exact-convolution}
\end{equation}
If $\nu/(r_*\gamma d_*)\le1/2$, then
\begin{equation}
 \sup_{z\in\mathcal D}|a_{\gamma,\nu}(z)|\le A_{\rm den}B.
 \label{eq:psi-density-bound}
\end{equation}
For instance, $0<\eta\le r_*\lambda d_*/4$ ensures this degree condition.
\end{lemma}
\begin{proof}
Taylor expansion of any polynomial $a$ terminates, so evenness of $K$ gives
\[
 \int_\R K(u)a(x-u/\gamma)\,du=M_K(D/\gamma)a(x),\qquad D=d/dx.
\]
All integrals exist by the moment hypothesis. Since $D$ is nilpotent on
polynomials of degree at most $\nu$, the finite substitution of the inverse
series gives \qeqref{eq:psi-exact-convolution} exactly.
Cauchy's estimates give $|b_{2\ell}|\le C_{\rm inv}r_*^{-2\ell}$ and
$|p_\nu^{(2\ell)}(z)|\le A_{\rm pol}B(2\ell)!/d_*^{2\ell}$ for
$z\in\mathcal D$. Using $(2\ell)!\le\nu^{2\ell}$ and summing a geometric
series of ratio at most $1/4$ proves \qeqref{eq:psi-density-bound}.
\end{proof}
We will also need an exterior bound, since the convolution in
\qeqref{eq:psi-exact-convolution} is over the whole real line. The real
Chebyshev extremal inequality gives, for $|t|\ge L_*$,
\begin{equation}
 |a_{\gamma,\nu}(t)|
 \le A_{\rm den}B\exp\left(\nu\operatorname{arcosh}\frac{|t|}{L_*}\right)
 \le A_{\rm den}B e^{\nu C_*(|t|-1)}.
 \label{eq:psi-exterior}
\end{equation}
The extremal inequality follows from alternation at the $\nu+1$ Chebyshev
extrema: a larger value outside $J_*$ would force the difference from a
suitable multiple of $T_\nu(t/L_*)$ to have more than $\nu$ zeros.
The last inequality uses $\operatorname{arcosh}s\le\log(2s)\le2s$ and
$|t|/L_*\le(|t|-1)/(L_*-1)$.

\subsubsection{Kernel identities and admissibility}
Let
\[
 \varphi(z)=\frac{e^{-z^2/2}}{\sqrt{2\pi}},\qquad
 \Phi(z)=\tfrac12+\int_0^z\varphi(t)\,dt,
 \qquad \varsigma(z)=(1+e^{-z})^{-1},
\]
\[
 \psi_G(z)=z\Phi(z),\qquad \psi_S(z)=z\varsigma(z).
\]
These are exact GELU and SiLU. Differentiation gives
\begin{align}
 K_G(z)=\psi_G''(z)&=(2-z^2)\frac{e^{-z^2/2}}{\sqrt{2\pi}},
 &M_G(z)&=(1-z^2)e^{z^2/2}, \label{eq:gelu-kernel}\\
 K_S(z)=\psi_S''(z)&=2g(z)+zg'(z),\quad g=\varsigma'=\tfrac14\sech^2(z/2),
 &M_S(z)&=\frac{(\pi z)^2\cos(\pi z)}{\sin^2(\pi z)}.
 \label{eq:silu-kernel}
\end{align}
At zero, both moment functions have removable value one. For GELU the moment
identity follows by differentiating $\int\varphi(t)e^{zt}\,dt=e^{z^2/2}$
twice. For SiLU, the substitution $v=e^t$ gives
\[
 G(z):=\int_\R g(t)e^{zt}\,dt
 =\int_0^\infty\frac{v^z}{(1+v)^2}\,dv
 =\Gamma(1+z)\Gamma(1-z)=\frac{\pi z}{\sin\pi z}
 \quad(|\Re z|<1).
\]
Integration by parts in $\int t g'(t)e^{zt}\,dt$ yields
$M_S=G-zG'$, proving \qeqref{eq:silu-kernel}.
Thus
\begin{equation}
 M_G^{-1}(z)=\frac{e^{-z^2/2}}{1-z^2},\qquad
 M_S^{-1}(z)=\frac{\sin^2(\pi z)}{(\pi z)^2\cos(\pi z)}.
 \label{eq:psi-inverse-symbols}
\end{equation}
The inverse radii are $1$ and $1/2$, respectively. For $0<r_*<1$ in the
GELU case and $0<r_*<1/2$ in the SiLU case, valid bounds are
\begin{equation}
 C_{{\rm inv},G}=\frac{e^{r_*^2/2}}{1-r_*^2},\qquad
 C_{{\rm inv},S}=\frac{[\sinh(\pi r_*)/(\pi r_*)]^2}{\cos(\pi r_*)}.
 \label{eq:psi-inverse-bounds}
\end{equation}
These follow from the power series of sine and
$|\cos(\pi z)|^2=\cos^2(\pi\Re z)+\sinh^2(\pi\Im z)$.

Choose $a_G>0$, $0<a_S<d_{\rm ext}/2$, and $0<\vartheta<\pi$, and set
\begin{equation}
 R_G=\lceil a_G\sqrt N\rceil,\quad R_S=\lceil a_SN\rceil,
 \quad r_R=(R+\tfrac12)h,\quad
 \sigma_G=\frac{2\pi h}{\lambda^2},\quad \sigma_S=\frac\vartheta\gamma.
 \label{eq:psi-geometry}
\end{equation}
For each activation choose $\eta>0$ as in Lemma~\ref{lem:psi-deconv}, with
$\eta\le1$. In the SiLU case also require $\eta\le\lambda/(8C_*)$.
In this subsection take $N$ large enough that $\nu\ge1$, $\gamma\ge1$, and
\begin{equation}
 r_R\le d_{\rm ext},\qquad \sigma_\psi\le s_*,\qquad
 \begin{cases}
 \nu C_*\le\gamma^2d_{\rm ext}/4,&\psi=\psi_G,\\
 \nu C_*\le\gamma/4,&\psi=\psi_S.
 \end{cases}
 \label{eq:psi-admissible}
\end{equation}
All these conditions hold for sufficiently large $N$ at fixed design. In
particular, the midpoint rectangle is contained in $\mathcal D$. The SiLU
contour stays below its nearest poles at preactivation height $\pi$.

\subsubsection{Halo and midpoint estimates}
Let $J_R=[a_R,b_R]=[-1-r_R,1+r_R]$ and write
\begin{align}
 (\mathcal C_Rp_\nu)(x)&=\int_{J_R}a_{\gamma,\nu}(t)K_\gamma(x-t)\,dt,\nonumber\\
 E_{\rm halo}^\psi(x)&=\int_{\R\setminus J_R}a_{\gamma,\nu}(t)K_\gamma(x-t)\,dt,
 \quad F_x^\psi(z)=a_{\gamma,\nu}(z)K_\gamma(x-z),\nonumber\\
 G_N^\psi(x)&=h\sum_{j=-R}^{N+R}a_{\gamma,\nu}(x_j)K_\gamma(x-x_j).
 \label{eq:psi-kernel-network}
\end{align}
There are no kernel poles in the chosen rectangles. Hence
Lemma~\ref{lem:midpoint} and \qeqref{eq:psi-exact-convolution} give
\begin{equation}
 f''-G_N^\psi=(f''-p_\nu)+E_{\rm halo}^\psi
             -\mathcal H_h[F_x^\psi]-\mathcal V_h[F_x^\psi].
 \label{eq:psi-master-decomp}
\end{equation}
Put $Z=\gamma r_R=\lambda(R+1/2)$. The following estimates include explicit
prefactors for the terms that will enter the approximation bound.

\begin{proposition}[GELU kernel quadrature]
\label{prop:gelu-errors}
Define
\begin{align}
 Q_G(\lambda)&=\frac{2(3+4\pi^2/\lambda^2)e^{-2\pi^2/\lambda^2}}
                    {1-e^{-4\pi^2/\lambda^2}},\nonumber\\
 V_G(\lambda)&=\frac4{\sqrt{2\pi}}
                 \left(\frac{4\lambda}\pi+\frac{2\lambda^3}{\pi^3}\right).
 \label{eq:gelu-allowances}
\end{align}
Under \qeqref{eq:psi-admissible}, uniformly on $\Om$,
\begin{equation}
 |E_{\rm halo}^G|\le12A_{\rm den}B e^{-Z^2/8},\quad
 |\mathcal H_h[F_x^G]|\le A_{\rm den}B Q_G(\lambda),\quad
 |\mathcal V_h[F_x^G]|\le A_{\rm den}B V_G(\lambda)e^{-Z^2/4}.
 \label{eq:gelu-errors}
\end{equation}
\end{proposition}
\begin{proof}
For real $u$,
$|K_G(u)|\le(2+u^2)e^{-u^2/2}/\sqrt{2\pi}$.
On $J_*\setminus J_R$, use the bounded density. Outside $J_*$, let
$w=|t-x|\ge d_{\rm ext}$. Equations \qeqref{eq:psi-exterior} and
\qeqref{eq:psi-admissible} imply
$|a_{\gamma,\nu}(t)|\le A_{\rm den}B e^{\gamma^2w^2/4}$.
Thus on both omitted tails, with $u=\gamma w$, the absolute integrand is at
most $A_{\rm den}B\gamma(2+u^2)e^{-u^2/4}/\sqrt{2\pi}$. Integration gives
\[
 |E_{\rm halo}^G|\le\frac{2A_{\rm den}B}{\sqrt{2\pi}}
 \int_Z^\infty(2+u^2)e^{-u^2/4}\,du
 \le12A_{\rm den}B e^{-Z^2/8},
\]
using $\int_0^\infty(2+u^2)e^{-u^2/8}\,du=6\sqrt{2\pi}$.

For the horizontal sides, for any real $v$,
\[
 \int_\R |K_G(u+iv)|\,du\le(3+v^2)e^{v^2/2}.
\]
Choose $v_G=\gamma\sigma_G=2\pi/\lambda$ and use \qeqref{eq:Rbounds}.
The two horizontal sides contribute at most
$2A_{\rm den}B(3+v_G^2)e^{v_G^2/2}/(e^{2\pi v_G/\lambda}-1)$,
which is $A_{\rm den}B Q_G$.

On a vertical side, the real kernel argument obeys $|u|\ge Z$.
The elementary inequality
$(2+u^2+v^2)e^{-u^2/2}\le(4+v^2)e^{-Z^2/4}$ gives
\[
 |\mathcal V_h[F_x^G]|\le
 \frac{4A_{\rm den}B}{\sqrt{2\pi}}e^{-Z^2/4}
 \int_0^{v_G}\frac{(4+v^2)e^{v^2/2}}{1+e^{2\pi v/\lambda}}\,dv.
\]
For $0\le v\le v_G$, the exponential ratio is at most $e^{-\pi v/\lambda}$.
Extend the integral to infinity to obtain $V_G$.
\end{proof}

\begin{proposition}[SiLU kernel quadrature]
\label{prop:silu-errors}
Let $c_\vartheta=\sec(\vartheta/2)$ and define
\begin{equation}
 A_\vartheta=2c_\vartheta^2+(1+\vartheta)c_\vartheta^3,
 \quad Q_S(\lambda)=\frac{8A_\vartheta}{e^{2\pi\vartheta/\lambda}-1},
 \quad V_S(\lambda)=\frac{4\lambda A_\vartheta\log2}\pi.
 \label{eq:silu-allowances}
\end{equation}
Under \qeqref{eq:psi-admissible}, uniformly on $\Om$,
\begin{equation}
 |E_{\rm halo}^S|\le16A_{\rm den}B e^{-Z/2},\quad
 |\mathcal H_h[F_x^S]|\le A_{\rm den}B Q_S(\lambda),\quad
 |\mathcal V_h[F_x^S]|\le A_{\rm den}B V_S(\lambda)e^{-Z/2}.
 \label{eq:silu-errors}
\end{equation}
\end{proposition}
\begin{proof}
On the real axis $|g(u)|\le e^{-|u|}$ and
$g'=g(1-2\varsigma)$, so $|K_S(u)|\le(2+|u|)e^{-|u|}$.
Outside $J_*$, \qeqref{eq:psi-exterior} and
$\nu C_*\le\gamma/4$ bound the density by
$A_{\rm den}B e^{\gamma|t-x|/4}$. Within $J_*$ use its uniform bound.
Consequently
\[
 |E_{\rm halo}^S|\le2A_{\rm den}B
 \int_Z^\infty(2+u)e^{-3u/4}\,du
 \le16A_{\rm den}B e^{-Z/2},
\]
since $(2+u)e^{-u/4}\le4$ for $u\ge0$.
For $z=u+iv$, $|v|\le\vartheta$, writing $t=e^{-u}$ gives
\[
 |1+te^{-iv}|\ge(1+t)\cos(\vartheta/2),\quad
 |g(z)|\le c_\vartheta^2e^{-|u|},\quad
 |1-2\varsigma(z)|\le c_\vartheta.
\]
Therefore $|K_S(u+iv)|\le A_\vartheta(1+|u|)e^{-|u|}$, whose real-line
integral is at most $4A_\vartheta$. The two horizontal sides and
\qeqref{eq:Rbounds} give $Q_S$. On each vertical side $|u|\ge Z$, and
$(1+|u|)e^{-|u|}\le2e^{-|u|/2}$. Four endpoint terms and
$\int_0^\infty(1+e^{2\pi y/h})^{-1}dy=h\log2/(2\pi)$ give $V_S$.
\end{proof}

\subsubsection{Activation realization and coefficient envelopes}
Define
\begin{equation}
 q_0^\psi(x)=\frac h\gamma\sum_{j=-R}^{N+R}
 a_{\gamma,\nu}(x_j)\psi(\gamma(x-x_j)).
 \label{eq:psi-q0}
\end{equation}
Twice differentiating gives $(q_0^\psi)''=G_N^\psi$. Choose
\begin{equation}
 \beta=f'(-1)-(q_0^\psi)'(-1),\quad
 \alpha=f(-1)-q_0^\psi(-1)+\beta,\quad
 q_\psi^\star=q_0^\psi+\alpha+\beta x.
 \label{eq:psi-anchor}
\end{equation}
Both activations obey $\psi(z)-\psi(-z)=z$ on their domains, by
$\Phi(-z)=1-\Phi(z)$ or $\varsigma(-z)=1-\varsigma(z)$.
Thus
\begin{equation}
 \beta x=\frac\beta\gamma[\psi(\gamma x)-\psi(-\gamma x)]
 \label{eq:psi-linear-identity}
\end{equation}
uses two additional hidden units and $\alpha$ is the output bias. The two
anchor conditions imply
\begin{equation}
 f(x)-q_\psi^\star(x)=\int_{-1}^x(x-t)[f''(t)-G_N^\psi(t)]\,dt.
 \label{eq:psi-lift-error}
\end{equation}
Integration gives a factor at most two in $L^\infty$; Young's inequality
with $s\mathbf1_{[0,2]}(s)$ gives the same factor in normalized $L^2$.

\begin{theorem}[Exact-GELU and SiLU witnesses]
\label{thm:psi-witness}\label{thm:gelu-witness}\label{thm:silu-witness}
With the domains, parameters, and admissibility conditions above, the
prescribed network has $W_\psi=N+2R_\psi+3$ hidden-unit slots and satisfies
\begin{align}
 \norm{f-q_G^\star}_\infty
 &\le 2B\left[A_{\rm pol}\rho_2^{-\nu}
       +A_{{\rm den},G}\{(12+V_G)e^{-Z_G^2/8}+Q_G\}\right]=:E_W^G,
 \label{eq:gelu-witness}\\
 \norm{f-q_S^\star}_\infty
 &\le 2B\left[A_{\rm pol}\rho_2^{-\nu}
       +A_{{\rm den},S}\{(16+V_S)e^{-Z_S/2}+Q_S\}\right]=:E_W^S.
 \label{eq:silu-witness}
\end{align}
The same estimates hold in normalized $L^2$. For fixed design these imply
\begin{equation}
 \norm{f-q_\psi^\star}_\infty
 \le C_\psi B e^{-c_\psi W_\psi}+2A_{{\rm den},\psi}B Q_\psi(\lambda),
 \label{eq:psi-rate}
\end{equation}
where $W_G=N+O(\sqrt N)$, $W_S=\Theta(N)$,
$Q_G=O((1+\lambda^{-2})e^{-2\pi^2/\lambda^2})$, and, for fixed
$\vartheta$, $Q_S=O_\vartheta(e^{-2\pi\vartheta/\lambda})$ on $0<\lambda\le1$.
\end{theorem}
\begin{proof}
Use \qeqref{eq:psi-master-decomp}, the two quadrature propositions, and
\qeqref{eq:psi-lift-error}. For $N\ge2/\eta$, $\nu\ge\eta N/2$.
Also $Z_G^2\ge\lambda^2a_G^2N$, $Z_S\ge\lambda a_SN$, and
$W_\psi\le(L_*+2)N$. One may therefore take
\[
 c_G=\frac{\min\{(\eta/2)\log\rho_2,\lambda^2a_G^2/8\}}{L_*+2},\qquad
 c_S=\frac{\min\{(\eta/2)\log\rho_2,\lambda a_S/2\}}{L_*+2}.
\]
The stated bounds on $Q_\psi$ follow from their explicit denominators.
\end{proof}
Unlike the tanh correction, the GELU halo estimate uses Gaussian spatial
decay directly. The uncorrected SiLU kernel has exponential tails and uses
$R_S=\Theta(N)$ to retain a geometric rate in total width.

On the real axis, $|\psi(t)|\le|t|$ and $|\psi'(t)|\le2$ for both
activations: $\psi_G'=\Phi+t\varphi$ has absolute value at most
$1+1/\sqrt{2\pi e}$, and $\psi_S'=\varsigma+t g$ has absolute value at most
$1+1/e$. Since $(N+2R+1)h\le2L_*$, write $H_*=1+L_*$ and set
\begin{equation}
 A_\beta=\delta^{-1}+4L_*A_{\rm den},\qquad
 A_\alpha=1+2L_*H_*A_{\rm den}+A_\beta.
 \label{eq:psi-affine-envelopes}
\end{equation}
Then \qeqref{eq:psi-density-bound} and \qeqref{eq:psi-anchor} give
\begin{equation}
 |w_j^\star|\le B A_{\rm den}h/\gamma,\quad
 |w_\pm^\star|=|\beta|/\gamma\le BA_\beta/\gamma,\quad
 |\alpha|\le BA_\alpha.
 \label{eq:psi-weight-envelope}
\end{equation}
Here $w_\pm^\star$ are the weights of the two affine units. These inequalities
follow respectively from the density bound,
$|(q_0^\psi)'(-1)|\le4L_*A_{\rm den}B$, and
$|q_0^\psi(-1)|\le2L_*H_*A_{\rm den}B$.

Let $\phi_j$ denote each native activation feature, including the affine
pair, and $\phi_b=1$. Choose deterministic positive normalization constants
$m_b=1$ and
\begin{equation}
 m_j\ge\max\{1,\sup_{x\in\Om}|\phi_j(x)|\},\qquad
 m_j\le2(1+\gamma H_*)
 \label{eq:psi-column-normalization}
\end{equation}
for the grid units, and $m_\pm\le2(1+\gamma)$ for the affine pair.
The exact supremum satisfies these conditions; the least dyadic upper bound on
$1+\gamma H_*$ (and $1+\gamma$ for the pair) also does so. These choices
may be validated by outward enclosures of the displayed quantities. Write $\varphi_j=\phi_j/m_j$ and $v_j=m_jc_j$.
Explicit normalized-coordinate envelopes are
\begin{equation}
 \alpha_j=m_jA_{\rm den}h/\gamma,\quad
 \alpha_\pm=m_\pm A_\beta/\gamma,\quad
 \alpha_b=A_\alpha.
 \label{eq:psi-coordinate-envelopes}
\end{equation}
Since $\gamma\ge1$, their sum satisfies
\begin{equation}
 \sum_j\alpha_j\le
 A_\alpha+4L_*A_{\rm den}(H_*+1)+8A_\beta=:K_{{\rm ref},\psi}.
 \label{eq:psi-normalized-l1}
\end{equation}
In particular, this bound includes the bias and both affine units.

\begin{lemma}[Normalized features on cell ellipses]
\label{lem:psi-strip-growth}
For every fixed $v_0>0$ for GELU, and $0<v_0<\pi$ for SiLU, there is a
constant $C_{\psi,v_0}$ such that
\begin{equation}
 |\psi(u+iv)|\le C_{\psi,v_0}(1+u_+),\qquad |v|\le v_0,
 \quad u_+=\max\{u,0\}.
 \label{eq:psi-strip-growth}
\end{equation}
Explicit choices are
\begin{equation}
 C_{G,v}=1+v+\frac{v e^{v^2/2}}{\sqrt{2\pi}}(e^{-1/2}+v),
 \qquad C_{S,v}=(1+v)\sec(v/2).
 \label{eq:psi-strip-constants}
\end{equation}
For a core cell of length $h$ and a Bernstein ellipse of parameter
$\varrho>1$, set $t_\varrho=(\varrho+\varrho^{-1}-2)/4$.
If $\lambda(\varrho-\varrho^{-1})/4\le v_0$, then every normalized feature
satisfies
\begin{equation}
 |\varphi_j(z)|\le C_{\psi,v_0}(3+\lambda t_\varrho)
 \quad\text{on that ellipse}.
 \label{eq:psi-ellipse-bound}
\end{equation}
\end{lemma}
\begin{proof}
For GELU, $|u|\Phi(u)$ is bounded for $u\le0$ by the Gaussian tail
inequality $t\Phi(-t)\le\varphi(t)$. In
\[
 \psi_G(u+iv)=(u+iv)\Phi(u)
        +i(u+iv)\int_0^v\varphi(u+it)\,dt,
\]
the second term is bounded uniformly in $u$, because
$|\varphi(u+it)|=e^{-u^2/2+t^2/2}/\sqrt{2\pi}$ and
$(|u|+v_0)e^{-u^2/2}$ is bounded. This proves \qeqref{eq:psi-strip-growth}.
For SiLU,
$|\psi_S(u+iv)|\le\sec(v_0/2)(|u|+v_0)/(1+e^{-u})$;
for $u<0$ use $|u|e^{u}\le1/e$, and for $u\ge0$ use the linear bound.

Each real feature's maximum preactivation on $\Om$ is $U_j=\gamma(1-x_j)$,
or $U_j=\gamma$ for an affine unit. Both activations satisfy
$\psi(t)\ge t/2$ for $t\ge0$. Consequently
$m_j\ge\max\{1,(U_j)_+/2\}$. The ellipse extends horizontally by at most
$h t_\varrho$ beyond the core; its preactivation real part is at most
$U_j+\lambda t_\varrho$. Divide \qeqref{eq:psi-strip-growth} by $m_j$ to
obtain \qeqref{eq:psi-ellipse-bound}.
\end{proof}

\paragraph{Recovery and computed error.}
Theorem~\ref{thm:psi-witness} supplies $E_W^{\rm an}=E_W^\psi$;
\qeqref{eq:psi-column-normalization}--\qeqref{eq:psi-normalized-l1} supply
the real-axis normalization and deterministic envelopes, including the bias
and both affine units. Under \qeqref{eq:sampling-conditions},
Lemma~\ref{lem:psi-strip-growth} supplies
$C_{\rm ell}=C_{\psi,\pi/4}(3+\lambda t_\varrho)$, with the SiLU ellipses
inside its pole-free strip. Thus the common dictionary assumptions in
Section~\ref{sec:recovery} hold. Under that section's numerical,
input-rounding, and evaluation conditions, with $M\ge n$, the same
computation proves
\begin{equation}
 \norm{f-\widehat q_W^\psi}_{2,\Om}
 \le4E_W^\psi+\eps_{\rm rec,W}^\psi,
 \qquad
 \norm{\widehat c_W^\psi}_1\le2BK_{{\rm ref},\psi}+e_{c,W}^\psi.
 \label{eq:psi-computed}
\end{equation}
In particular, \qeqref{eq:psi-rate} yields geometric refinement in width plus
the activation-specific midpoint and numerical allowances.
Corollary~\ref{cor:logwidth} applies under its fixed-design and certified-range
premises.

